\documentclass{beamer}
\usepackage{listings}
\usepackage{color}

\definecolor{codegreen}{rgb}{0,0.6,0}
\definecolor{codegray}{rgb}{0.5,0.5,0.5}
\definecolor{codepurple}{rgb}{0.58,0,0.82}
\definecolor{backcolour}{rgb}{0.95,0.95,0.92}

\lstdefinestyle{mystyle}{
    backgroundcolor=\color{backcolour},   
    commentstyle=\color{codegreen},
    keywordstyle=\color{magenta},
    numberstyle=\tiny\color{codegray},
    stringstyle=\color{codepurple},
    basicstyle=\footnotesize,
    breakatwhitespace=false,         
    breaklines=true,                 
    captionpos=b,                    
    keepspaces=true,                 
    numbers=left,                    
    numbersep=5pt,                  
    showspaces=false,                
    showstringspaces=false,
    showtabs=false,                  
    tabsize=2
}

\lstset{style=mystyle}

\title{Static vs. Dynamic Scoping}
\author{Your Name}
\date{\today}

\begin{document}

\begin{frame}
  \titlepage
\end{frame}

\begin{frame}{Introduction to Scoping}
  \begin{itemize}
    \item Scoping determines how variables are accessed.
    \item It's like rules for finding the "right box" (variable).
    \item Two main types: Static and Dynamic.
  \end{itemize}
\end{frame}

\begin{frame}[fragile]{Static Scoping (Python)}
  \begin{itemize}
    \item Scope determined by code structure (lexical).
    \item Determined at compile time.
  \end{itemize}
  \begin{lstlisting}[language=Python]
x = 10

def outer():
    x = 20
    def inner():
        print("Inside Inner (Python):", x)
    inner()
    standalone()

def standalone():
    print("Inside Standalone (Python):", x)

outer()
standalone()
  \end{lstlisting}
\end{frame}

% \begin{frame}{Static Scoping Explanation}
%   \begin{itemize}
%     \item \texttt{inner()} finds \texttt{x = 20} in \texttt{outer()}'s scope.
%     \item \texttt{standalone()} finds \texttt{x = 10} in the global scope.
%     \item Output:
%       \begin{verbatim}
% Inside Inner (Python): 20
% Inside Standalone (Python): 10
% Inside Standalone (Python): 10
%       \end{verbatim}
%   \end{itemize}
% \end{frame}

\begin{frame}[fragile]{Dynamic Scoping (Bash)}
  \begin{itemize}
    \item Scope determined by call stack (runtime).
    \item Depends on the calling sequence.
  \end{itemize}
  \begin{lstlisting}[language=bash]
#!/bin/bash

x=10

function outer {
  local x=20
  inner
  standalone
}

function inner {
  echo "Inside Inner (Bash): $x"
}

function standalone {
  echo "Inside Standalone (Bash): $x"
}

outer
  \end{lstlisting}
\end{frame}

% \begin{frame}{Dynamic Scoping Explanation}
%   \begin{itemize}
%     \item \texttt{inner()} and \texttt{standalone()} find \texttt{x = 20} in \texttt{outer()}'s scope.
%     \item Output:
%       \begin{verbatim}
% Inside Inner (Bash): 20
% Inside Standalone (Bash): 20
%       \end{verbatim}
%   \end{itemize}
% \end{frame}

% \begin{frame}{Key Differences}
%   \begin{itemize}
%     \item \textbf{Python (Static):} Definition location matters.
%     \item \textbf{Bash (Dynamic):} Call location matters.
%     \item Static: Predictable, easier to debug.
%     \item Dynamic: Runtime customizations, but harder to reason about.
%   \end{itemize}
% \end{frame}

% \begin{frame}{Summary}
%     \begin{itemize}
%         \item Static scoping uses the location of the code to determine the scope.
%         \item Dynamic scoping uses the calling sequence to determine scope.
%         \item Python uses Static, Bash uses Dynamic.
%     \end{itemize}
% \end{frame}

\end{document}